\subsection{Audit of every use of the inverse}

The singular matrix $D$ from the original problem is never inverted. It is
first split into active and inactive blocks, and the inactive block
disappears from $ADA^{\mathsf T}$. The only inverse of a diagonal weight
matrix in the proof is $D_{+,d}^{-1}$ inside the active weighted-Wishart
linearization, where
\begin{equation}
\|D_{+,d}^{-1}\|\le\tau^{-1}.
\end{equation}
The exact Schur identities of Lemma~\ref{S:exact-schur-lemma} themselves do
not require a weight inverse; they require only that the chosen spectral
parameter avoid the orthogonal block. No Moore--Penrose inverse is used, no
zero mode is removed from the original matrix, and the exact nullity is
proved separately below.

\section{Random active size}

The marked theorem is deterministic in the active sample size. The following
lemma transfers all of its high-probability conclusions to the binomially
random active size without requiring a theorem uniform over all aspect ratios.

\begin{lemma}[Random-index transfer]
\label{S:random-index-lemma}
For each $d$ and integer $m$, let $r_{d,m}\in[0,1]$ be the conditional failure probability of a
specified active-model event under $m$ independent active marks. Suppose that
$r_{d,m_d}\to0$ for every deterministic sequence $m_d/d\to c$. If $N_d/d\to c$ in probability,
then
\begin{equation}
\E r_{d,N_d}\longrightarrow0.
\label{S:random-index}
\end{equation}
\end{lemma}

\begin{proof}
Choose $\eta_d\downarrow0$ so slowly that
$\Pp\{|N_d/d-c|>\eta_d\}\to0$. If
\begin{equation}
\sup_{\substack{m\in\mathbb N\\|m/d-c|\le\eta_d}}r_{d,m}
\end{equation}
did not tend to zero, there would be a subsequence and integers $m_d$ in the displayed window
with $m_d/d\to c$ but $r_{d,m_d}$ bounded away from zero, contradicting the hypothesis.
Splitting $\E r_{d,N_d}$ on the typical and atypical events proves \eqref{S:random-index}.
\end{proof}

